\documentclass[11pt,a4paper]{article}

\usepackage[utf8]{inputenc}
\usepackage[T1]{fontenc}
\usepackage{lmodern}
\usepackage[margin=2.3cm]{geometry}
\usepackage{amsmath,amssymb,bm}
\usepackage{mhchem}
\usepackage{siunitx}
\usepackage{graphicx}
\usepackage{booktabs,tabularx}
\usepackage{enumitem}
\usepackage[dvipsnames]{xcolor}
\usepackage{tcolorbox}
\usepackage{tikz}
\usetikzlibrary{arrows.meta,positioning,shapes.geometric}
\usepackage{listings}
\usepackage{microtype}
\usepackage[colorlinks=true,linkcolor=MidnightBlue,citecolor=ForestGreen,urlcolor=MidnightBlue]{hyperref}

\sisetup{detect-all}
\lstset{basicstyle=\ttfamily\small,columns=fullflexible,keepspaces=true,
        frame=single,framerule=0.3pt,rulecolor=\color{gray!60},backgroundcolor=\color{gray!6}}

\newtcolorbox{keybox}[1]{colback=MidnightBlue!4,colframe=MidnightBlue!70,
  fonttitle=\bfseries,title=#1,boxrule=0.6pt,arc=2pt}
\newtcolorbox{warnbox}[1]{colback=BrickRed!4,colframe=BrickRed!70,
  fonttitle=\bfseries,title=#1,boxrule=0.6pt,arc=2pt}
\newtcolorbox{nowbox}[1]{colback=ForestGreen!4,colframe=ForestGreen!70,
  fonttitle=\bfseries,title=#1,boxrule=0.6pt,arc=2pt}
\newcommand{\question}[1]{\par\smallskip\noindent\textcolor{MidnightBlue}{\textbf{Q.}~\emph{#1}}\par\smallskip}


\title{\textbf{A neural network potential for water clusters}\\[4pt]
\large Lesson notes accompanying \texttt{tutorial\_02\_wclusters/wclusters.ipynb}}
\author{Machine-learned interatomic potentials --- hands-on session}
\date{}

\begin{document}
\maketitle

\begin{abstract}
The second notebook applies the workflow of the Zundel tutorial --- committee
neural networks, query by committee (QbC), MD in CP2K, validation by
observables --- to a protonated water cluster \ce{(H2O)6H+}. It then asks the
question that matters most in practice: \emph{how far does the model
transfer?} It tests the model on a cluster almost three times larger,
\ce{(H2O)16H+}, and on periodic bulk water. Finally, it questions the
reference method itself. This is still a \textbf{basic, pedagogical
example}: small clusters in vacuum, a semi-empirical reference and a tiny
network. Its failures are as instructive as its successes. These notes assume
you have done the Zundel tutorial and its lesson. The last section extends its
discussion of the subtleties beyond the simple model with points specific to
water.
\end{abstract}

\tableofcontents

%======================================================================
\section{From one ion to a cluster}
%======================================================================

The Zundel ion had 7 atoms, all within a single cutoff sphere. Every atom
``saw'' the whole molecule, and the network effectively learned one global
function. \ce{(H2O)6H+} (19 atoms) is the first system where the
\emph{locality} of the Behler--Parrinello model really matters:
\begin{itemize}[nosep]
  \item The cluster is $\sim\SI{7.5}{\angstrom}$ across, more than the cutoff
        $r_c = 12\,a_0 \approx \SI{6.35}{\angstrom}$, so atoms on opposite sides do
        not see each other.
  \item It contains a hydrogen-bond network. The excess proton can take
        Zundel-like (\ce{H5O2+}) and Eigen-like (\ce{H9O4+}) local
        configurations and hop between waters~\cite{Marx1999}. The model must
        describe both, from local information only.
  \item There are now many oxygens, so the O--O--O angular functions excluded
        in tutorial~1 are switched back on (\texttt{exclude\_triples} is
        commented out). Both elements now have 32 descriptors and 201 weights.
\end{itemize}
Protonated water clusters are a classic test case for MLIPs. They have been
fitted up to coupled-cluster accuracy with exactly this type of model and an
automated, committee-based selection of training
data~\cite{Schran2020cc}.

\begin{keybox}{What this example is --- and what it is not}
\begin{tabularx}{\linewidth}{@{}lX@{}}
\textbf{Training system} & \ce{(H2O)6H+}, 19 atoms, vacuum.\\
\textbf{Test systems} & \ce{(H2O)16H+} (49 atoms, vacuum); 64 \ce{H2O} in a periodic box (192 atoms, $NpT$).\\
\textbf{Reference} & SCC-DFTB~\cite{Elstner1998}. Known to describe water--water interactions poorly (Sec.~\ref{sec:bulk}).\\
\textbf{Model} & Same as tutorial 1: 2 hidden layers $\times$ 5 neurons per element, 8-member committee.\\
\end{tabularx}
\medskip
A research model for aqueous systems would use a much better reference
(hybrid DFT with dispersion, or coupled cluster), bigger or equivariant
networks, and above all training data that include the condensed phase
(Sec.~\ref{sec:novelty}).
\end{keybox}

%======================================================================
\section{Reference data}
%======================================================================

All reference trajectories in \texttt{data/} were generated with SCC-DFTB in
CP2K, with a Nos\'e--Hoover chain thermostat at \SI{300}{K}
($\tau=\SI{50}{fs}$) and $\Delta t = \SI{0.5}{fs}$:
\begin{center}\small
\begin{tabular}{@{}lllll@{}}
\toprule
directory & system & atoms & ensemble, length & role\\
\midrule
\texttt{data/cluster}     & \ce{(H2O)6H+}  & 19  & NVT, \SI{25}{ps}, 10\,001 frames & training + validation\\
\texttt{data/big\_cluster} & \ce{(H2O)16H+} & 49  & NVT, \SI{25}{ps}, 10\,001 frames & validation only\\
\texttt{data/bulk}        & 64 \ce{H2O}    & 192 & $NpT$, \SI{1}{atm}, \SI{20}{ps} & density only (cell file)\\
\bottomrule
\end{tabular}
\end{center}
The key point: \textbf{the model is trained only on the small cluster.}
Everything else tests transferability.

\question{Before running anything: which local environments occur in bulk
water but never in a 7-water cluster? Think about the number of neighbours
inside a \SI{6.35}{\angstrom} sphere, and about surface molecules.}

%======================================================================
\section{Active learning on the cluster}
%======================================================================

This notebook skips the naive single model and goes directly to QbC, with the
same settings as tutorial~1: 20 random initial structures, then 15 iterations
that each add the 20 structures (out of 500 random candidates) with the largest
force disagreement, using 8 members and 15 epochs. The final committee is
retrained for 50 epochs on the last training set.

\begin{table}[h]
\centering\small
\caption{Mean committee force disagreement (\si{meV/\angstrom}) during QbC on
\ce{(H2O)6H+} in a test run, and wall time per iteration (8 cores).}
\label{tab:qbc}
\begin{tabular}{@{}rrrrrr@{}}
\toprule
iteration & training-set size & train & candidates & selected & time (s)\\
\midrule
1  & 20  & 129 & 151 & 228 & 12\\
2  & 40  & 100 & 102 & 142 & 13\\
3  & 60  &  79 &  75 &  98 & 13\\
5  & 100 &  62 &  56 &  73 & 14\\
10 & 193 &  50 &  43 &  51 & 16\\
15 & 281 &  45 &  38 &  49 & 19\\
\bottomrule
\end{tabular}
\end{table}

Compare with the Zundel run (tutorial~1 lesson, Table~1):
\begin{itemize}[nosep]
  \item The initial disagreement is much lower (151 vs.\ \SI{695}{meV/\angstrom}
        on candidates). The cluster offers many more atoms per structure, and
        hence many more force labels, for the same 20 structures.
  \item The final disagreement is similar ($\sim\SI{40}{meV/\angstrom}$), with a
        similar number of training structures ($\sim$280).
  \item This time the cost per iteration \emph{grows}, from 12 to \SI{19}{s}.
        Training cost scales with the number of atoms in the training set, and
        each structure now has 19 atoms instead of 7.
  \item Late iterations re-select a few structures already in the set
        (``$n$ already present''), a sign that the pool is exhausted.
\end{itemize}

\question{The disagreement converges to $\sim\SI{40}{meV/\angstrom}$, not to
zero. What limits it: the data, the network size (10 hidden neurons), or the
locality approximation? How could you tell them apart?}

%======================================================================
\section{MD and validation on the training system}
%======================================================================

The committee runs \SI{10}{ps} of NVT MD on \ce{(H2O)6H+} (\texttt{nnp\_md/}).
In the test run it stayed stable: mean temperature \SI{297}{K} (reference
\SI{303}{K}), and no molecule left the cluster.

\begin{center}\small
\begin{tabular}{@{}lccc@{}}
\toprule
score (\%) & \ce{H5O2+} (tut.~1) & \ce{(H2O)6H+} & \ce{(H2O)16H+}\\
\midrule
RDF O--O & 92.8 & 87.4 & 66.3\\
RDF O--H & 97.0 & 91.0 & 82.9\\
RDF H--H & 97.5 & 94.8 & 77.5\\
VDOS (mean of O, H) & 90.6 & 79.9 & 82.4$^\ast$\\
\bottomrule
\end{tabular}\\[2pt]
{\footnotesize Score $=100(1-\varepsilon)$, with $\varepsilon$ the normalised
absolute deviation from the DFTB curve (tutorial~1 lesson, Sec.~6).
$^\ast$Recomputed; see the warning box in Sec.~\ref{sec:big}.}
\end{center}

The scores are lower than for the Zundel ion, even on the training system.
There are three reasons. The system is more complex, but the network is the
same size. The NNP run is shorter than the reference (10 vs.\ \SI{25}{ps}),
which adds statistical noise to the RDFs and, above all, to the spectra.
And a cluster in vacuum is a ``soft'' object: its shape fluctuates and
molecules can drift outwards.

\begin{warnbox}{Molecules flying away}
The notebook warns that you may see a water molecule leave the cluster, and
attributes it to DFTB, which \emph{significantly} underestimates water--water
interactions. This is not an NNP artefact: the DFTB reference does it too
(see next section). An MLIP reproduces the physics of its reference,
including the bad parts.
\end{warnbox}

%======================================================================
\section{Transfer 1: a larger cluster, \ce{(H2O)16H+}}
\label{sec:big}
%======================================================================

The same model now runs \SI{10}{ps} on a 49-atom cluster it has never seen
(\texttt{big\_cluster/}), and is compared with a \SI{25}{ps} DFTB reference.
The run is stable (mean \SI{302}{K}), but the RDF scores drop sharply,
especially O--O (\SI{66}{\%}).

\paragraph{Read the numbers critically.} In the test run:
\begin{itemize}[nosep]
  \item The \emph{DFTB reference} cluster loses water molecules. After
        \SI{11.8}{ps} one oxygen sits more than \SI{8}{\angstrom} from the
        centre of mass, and by \SI{25}{ps} three molecules have evaporated
        (up to \SI{15.6}{\angstrom} away).
  \item The NNP trajectory stops at \SI{10}{ps}, \emph{before} evaporation
        would have happened in the reference. One molecule made a temporary
        excursion beyond \SI{8}{\angstrom} at \SI{2.2}{ps} and came back.
\end{itemize}
Part of the low O--O score therefore comes from comparing \emph{different
physical situations}: an intact cluster against one that is falling
apart. It is not purely a model error. Comparisons by observables are only
meaningful for equivalent, converged simulations: same length, same
ensemble, same state of the system.

\begin{warnbox}{Bug in the notebook: wrong VDOS for the big cluster}
The last line of the big-cluster scoring cell reads
\begin{lstlisting}
mlps.run_vdos_test(ref_atv, dt_trj, nnp_atv, dt_trj)
\end{lstlisting}
It reuses the velocities of the \emph{small} cluster, so the printed VDOS
score is just a copy of the previous one (79.8 / 80.0\,\%). The correct call is
\begin{lstlisting}
mlps.run_vdos_test(ref_big_atv, dt_trj, nnp_big_atv, dt_trj)
\end{lstlisting}
which in the test run gave 87.3\,\% (O) and 77.5\,\% (H), a mean of 82.4\,\%.
\end{warnbox}

\question{(Notebook) The scores seem almost reasonable, but from a visual
inspection of the trajectories, are the dynamics well reproduced? What do you
conclude about accuracy criteria, and about transferability?}

%======================================================================
\section{Transfer 2: bulk water}
\label{sec:bulk}
%======================================================================

The last test is the hardest: 64 neutral water molecules in a periodic box,
$NpT$ at \SI{1}{atm} and \SI{300}{K}, starting at
$\rho = \SI{1.0}{g/cm^3}$ (\texttt{bulk/}, \SI{2.5}{ps}). Note the two
changes at once: no excess proton, and no vacuum.

\paragraph{What happened in the test run.} The NNP liquid \emph{expands}. The
density fell from 1.00 to \SI{0.52}{g/cm^3} within \SI{2.5}{ps}, while the
temperature peaked at \SI{1127}{K}. The model does not hold the liquid
together. In bulk, each molecule has neighbours in every direction inside the
cutoff, an environment entirely absent from a 7-water cluster, where every
molecule sits at a surface. Concretely, an oxygen in the 7-water cluster has
at most 6 other oxygens inside the cutoff; in the bulk box it has 28--38 (34 on
average). The radial descriptors, which essentially count neighbours, take
values far outside the training range. The network extrapolates, and the
predicted forces (and hence the pressure) are meaningless.

\paragraph{And the reference?} The notebook then plots the density of the
DFTB bulk simulation (\texttt{data/bulk/bulk-1.cell}). It goes the other
way: the box shrinks from 12.42 to \SI{11.13}{\angstrom}, and the density
settles around \SI{1.41}{g/cm^3}, about 40\,\% above the experimental
\SI{0.997}{g/cm^3}. This DFTB parametrisation produces a strongly
over-structured, over-dense liquid. Even a perfect MLIP trained on it would be
a perfect model of the wrong water.

\question{(Notebook) What is happening in the NNP bulk run? Keep in mind that
CP2K writes unwrapped coordinates, and that nglview does not handle periodic
boundary conditions, so molecules seem to ``leave'' the box even in a
healthy simulation. Why is the density (cell file) a more reliable
diagnostic here than the movie?}

\begin{keybox}{Lessons of this tutorial}
\begin{enumerate}[nosep,leftmargin=*]
  \item \textbf{Locality enables transfer, data limit it.} A local model can
        run on any system size, but it is only reliable for local environments
        represented in its training set.
  \item \textbf{Validate with care.} Observable-based scores depend on how the
        comparison is set up (run length, state of the system), and the code
        that computes them can have bugs.
  \item \textbf{The reference sets the ceiling.} Check the reference method
        against experiment or a higher level of theory \emph{before} investing
        in the ML model.
\end{enumerate}
\end{keybox}

%======================================================================
\section{Going further: subtleties beyond the simple model}
\label{sec:novelty}
%======================================================================

The general refinements --- equivariant message-passing networks, foundation
models and fine-tuning, long-range electrostatics, energy-conserving
benchmarks --- are summarised in Sec.~8 of the Zundel lesson. Here are the
points specific to this tutorial.

\paragraph{Clusters \emph{can} teach the bulk, if done right.} This tutorial
shows a cluster-trained model failing in bulk. Yet one study did exactly
the opposite, on purpose: it trained a Behler--Parrinello network on
\emph{finite clusters} computed with CCSD(T), the ``gold standard'' of quantum
chemistry, which is unaffordable for periodic liquids. The model then ran
coupled-cluster-quality MD of bulk liquid water, including nuclear quantum
effects~\cite{Daru2022}. The difference: their clusters are large and chosen
so that, within the cutoff, they contain bulk-like environments, and the
selection is automated~\cite{Schran2020cc}. Transferability is a property of
the \emph{training set}, not of the model.

\paragraph{Better references are affordable.} Because MLIPs are cheap, the
reference can be expensive. Hybrid DFT with dispersion is the usual choice
for aqueous systems~\cite{Schran2021}, and coupled-cluster-level models of
water exist~\cite{Daru2022}. A production water model would not be built on
SCC-DFTB. In this tutorial it is a deliberate choice, to keep the data
cheap and to show what a poor reference does.

\paragraph{Active learning in the loop.} Here QbC only selects from an
\emph{existing} DFTB trajectory. In research workflows the committee runs the
MD itself. Configurations with high disagreement (for example, the first
frames of the bulk run above) are sent for new reference calculations and
added to the training set, iteratively, until the disagreement stays low in
production~\cite{Schran2020,Schran2021}. Tutorial~3 uses models built this
way for water at interfaces, in confinement and with ions.

\paragraph{Foundation models as a starting point.} A pretrained universal
model~\cite{Batatia2023,Wood2025} would not explode in bulk water the way this
model does: it has seen condensed phases. But its accuracy for water depends
on the reference functional of its training data, and charged clusters such
as \ce{(H2O)_nH+} need a model trained on charged systems (e.g.\ on
OMol25~\cite{Levine2025}). A natural refinement is to fine-tune such a model on a
few hundred structures selected by committee disagreement. That is the same
idea as QbC, applied to a much better starting point.

\paragraph{Long-range electrostatics and the excess charge.} A $+1$ cluster
in vacuum has a net charge, and its surface molecules have a long-range
electrostatic environment that a $\SI{6.35}{\angstrom}$ cutoff cannot see.
Charge-aware architectures~\cite{Ko2021,Cheng2024} address exactly this.

%======================================================================
\section{Exercises}
%======================================================================
\begin{enumerate}
  \item Plot the training points selected by QbC on top of $E_p(t)$. Compare
        with the Zundel case.
  \item Change the architecture (\texttt{n\_hidden}, \texttt{n\_neurons}). Does
        the final committee disagreement go below
        $\sim\SI{40}{meV/\angstrom}$? (Remember to delete the old
        \texttt{qbc/iteration-*}, \texttt{qbc/results.shelf.*} and
        \texttt{train/final\_training} directories first, as the notebook says.)
  \item Fix the VDOS line for the big cluster and compare the scores.
  \item Make the big-cluster comparison fair: score the NNP run against only
        the first \SI{10}{ps} of the reference (use the \texttt{stop} argument of
        \texttt{mlps.load\_with\_cell}). How much do the RDF scores change?
  \item Monitor the cluster's integrity, e.g.\ the largest distance of an
        oxygen from the centre of mass, in the reference and NNP runs of both
        clusters.
  \item Bulk water: in \texttt{bulk/bulk.inp}, the per-member energies are
        switched off (\texttt{\&ENERGIES OFF} in the \texttt{\&NNP/\&PRINT}
        section; the same holds for the cluster inputs). Switch them on as in
        tutorial~1's \texttt{zundel.inp} (\texttt{\&ENERGIES SILENT} with
        \texttt{\&EACH MD 1}), rerun, and compute
        the committee energy spread along the trajectory. When does it blow
        up compared with the density? How would an active-learning loop use
        this?
  \item Discuss: what training set would you need for a model of bulk water?
        Of the air--water interface?
\end{enumerate}

%======================================================================
\section{Running times}
%======================================================================
Measured on an 8-core workstation (one core per committee member):

\begin{center}\small
\begin{tabular}{@{}lr@{}}
\toprule
step & wall time\\
\midrule
QbC (15 iterations, 8 members) & \SI{3.9}{min}\\
final training (8 members, 281 structures, 50 epochs) & \SI{48}{s}\\
\ce{(H2O)6H+} MD (\num{20000} steps, \SI{10}{ps}) & \SI{46}{s}\\
\ce{(H2O)16H+} MD (\num{20000} steps, \SI{10}{ps}) & \SI{1.6}{min}\\
bulk water $NpT$ MD (\num{5000} steps, \SI{2.5}{ps}, 192 atoms) & \SI{6.1}{min}\\
\bottomrule
\end{tabular}
\end{center}
The bulk run is the longest step. Start it as soon as the final model is
trained, and do the big-cluster analysis while it runs.

%======================================================================
\begin{thebibliography}{99}
\small

\bibitem{Marx1999} D.~Marx, M.~E.~Tuckerman, J.~Hutter and M.~Parrinello,
  \emph{The nature of the hydrated excess proton in water},
  Nature \textbf{397}, 601 (1999).

\bibitem{Schran2020cc} C.~Schran, J.~Behler and D.~Marx,
  \emph{Automated fitting of neural network potentials at coupled cluster
  accuracy: protonated water clusters as testing ground},
  J.\ Chem.\ Theory Comput.\ \textbf{16}, 88 (2020); arXiv:1908.08734.

\bibitem{Elstner1998} M.~Elstner \emph{et al.},
  \emph{Self-consistent-charge density-functional tight-binding method for
  simulations of complex materials properties}, Phys.\ Rev.\ B \textbf{58},
  7260 (1998).

\bibitem{Schran2020} C.~Schran, K.~Brezina and O.~Marsalek,
  \emph{Committee neural network potentials control generalization errors and
  enable active learning}, J.\ Chem.\ Phys.\ \textbf{153}, 104105 (2020).

\bibitem{Schran2021} C.~Schran, F.~L.~Thiemann, P.~Rowe, E.~A.~M\"uller,
  O.~Marsalek and A.~Michaelides,
  \emph{Machine learning potentials for complex aqueous systems made simple},
  Proc.\ Natl.\ Acad.\ Sci.\ USA \textbf{118}, e2110077118 (2021);
  arXiv:2106.00048.

\bibitem{Daru2022} J.~Daru, H.~Forbert, J.~Behler and D.~Marx,
  \emph{Coupled cluster molecular dynamics of condensed phase systems enabled
  by machine learning potentials: liquid water benchmark},
  Phys.\ Rev.\ Lett.\ \textbf{129}, 226001 (2022).

\bibitem{Batatia2023} I.~Batatia \emph{et al.},
  \emph{A foundation model for atomistic materials chemistry},
  arXiv:2401.00096 (2023).

\bibitem{Wood2025} B.~M.~Wood \emph{et al.},
  \emph{UMA: A family of universal models for atoms}, arXiv:2506.23971 (2025).

\bibitem{Levine2025} D.~S.~Levine \emph{et al.},
  \emph{The Open Molecules 2025 (OMol25) dataset, evaluations, and models},
  arXiv:2505.08762 (2025).

\bibitem{Ko2021} T.~W.~Ko, J.~A.~Finkler, S.~Goedecker and J.~Behler,
  \emph{A fourth-generation high-dimensional neural network potential with
  accurate electrostatics including non-local charge transfer},
  Nat.\ Commun.\ \textbf{12}, 398 (2021).

\bibitem{Cheng2024} B.~Cheng,
  \emph{Latent Ewald summation for machine learning of long-range
  interactions}, arXiv:2408.15165 (2024).

\end{thebibliography}

\end{document}
